This paper investigates how software engineering teams can successfully build and sustain diversity across race, ethnicity, and gender within a large technology organization. Although prior research suggests that diverse teams can improve innovation, productivity, and software quality, many software engineering teams remain demographically homogeneous. The authors argue that while many organizations publicly commit to diversity goals, there is limited empirical evidence on the specific practices that actually help teams achieve and maintain diversity. To address this gap, the study examines the experiences and practices of highly diverse engineering teams at Google in order to understand what actions contribute to building inclusive and representative teams.

To explore this question, the researchers conducted 20 semi-structured interviews with engineers and managers from some of the most diverse engineering teams within the organization. Using qualitative thematic analysis, the study identifies practices related to recruiting, hiring, and creating inclusive work environments that help sustain diversity. The findings highlight two major themes: first, strategies used by managers to intentionally recruit diverse candidates, and second, the role of inclusive practices—particularly the concept of \textit{technical allyship} — in supporting engineers from underrepresented groups. These insights provide actionable recommendations for engineers, managers, and organizational leaders seeking to improve diversity and inclusion in software engineering teams.

\section{Motivation}
Software engineering teams play a central role in shaping the technologies used by millions of people, yet these teams remain largely demographically homogeneous across many organizations. When teams lack diversity, the perspectives that influence how software is designed, implemented, and evaluated can become limited. \textbf{Research has shown that diversity within collaborative engineering teams can produce tangible benefits. For example, differences in cognitive styles among team members have been associated with increased novelty in software design decisions, while gender and tenure diversity have been linked with improved team productivity and performance outcomes} \cite{pretorius2020intuition,vasilescu2015gender}. These findings suggest that diversity is not only a social or ethical concern but also a factor that can shape the technical outcomes and effectiveness of software development teams.

Beyond productivity benefits, diversity can also influence the social and organizational dynamics of software teams. \textbf{Studies have found that teams with greater representation of women tend to experience fewer organizational issues known as community smells, which refer to problematic communication and collaboration patterns within development communities} \cite{catolino2019gender}. Other research indicates that increased gender diversity can help reduce biases among team members and foster more inclusive work environments \cite{wang2019implicit}. These insights have motivated many large technology companies to publicly commit to improving diversity and representation in their engineering workforce, often through organizational diversity initiatives and hiring goals.

Despite these commitments, there remains a significant gap between the recognition of diversity's benefits and the availability of evidence-based practices for achieving it. \textit{Much of the existing literature focuses on the consequences of diversity rather than on the processes that help organizations build diverse teams in the first place. As a result, recommendations for improving diversity often rely on anecdotal accounts or high-level organizational policies rather than empirically grounded practices that engineering teams and managers can implement in their day-to-day work} \cite{kreitz2008best}. This lack of actionable knowledge makes it difficult for organizations to translate diversity goals into practical strategies.

Another motivation arises from the persistent challenges faced by developers from underrepresented groups in the software industry. Prior research has documented several structural and cultural barriers affecting these developers, including bias in evaluation processes, discrimination in professional environments, isolation within teams, and lower acceptance rates of technical contributions in collaborative projects \cite{nadri2022relationship,terrell2017gender}. These challenges can discourage participation and retention, further reinforcing the lack of diversity in software engineering. Understanding how certain teams successfully overcome these challenges and cultivate inclusive environments therefore becomes essential for identifying practices that can help organizations recruit, support, and retain a more diverse engineering workforce.

\subsection{Related Work}

Research on diversity in software engineering and organizational settings has grown substantially in recent years. Much of this literature examines the consequences of diversity within teams, particularly how demographic or cognitive differences influence team performance, innovation, and collaboration outcomes. For example, studies of software development teams have shown that diversity in gender, experience, and tenure can positively affect productivity and collaborative dynamics \cite{vasilescu2015gender}. Similarly, work on employee diversity in innovation-driven environments suggests that teams composed of individuals with varied perspectives and backgrounds are more likely to produce novel ideas and solutions \cite{ostergaard2011different}. These studies provide empirical support for the argument that diversity can enhance the effectiveness of collaborative technical work.

Beyond performance outcomes, researchers have also investigated how diversity shapes social interactions and organizational structures within software development communities. Prior work has documented that gender diversity can influence collaboration patterns and reduce negative organizational phenomena such as community smells, which represent problematic coordination and communication structures within development teams \cite{catolino2019gender}. Other studies have explored the presence of implicit bias in professional software development environments, demonstrating that stereotypes and unconscious biases can influence evaluation, collaboration, and decision-making processes among developers \cite{wang2019implicit}. These findings highlight that diversity not only affects technical outcomes but also interacts with social and cultural dynamics within engineering teams.

Another important line of research focuses on the experiences and challenges faced by developers from underrepresented groups in the software industry. Empirical studies have documented barriers such as discrimination, social isolation, and reduced recognition of technical contributions. For instance, investigations of contribution evaluation in open-source software communities reveal disparities in how contributions from developers of different genders or perceived identities are reviewed and accepted \cite{terrell2017gender}. Similarly, analyses of demographic perceptions in collaborative software projects have shown that factors such as race and ethnicity can influence how contributions are evaluated and perceived by other developers \cite{nadri2022relationship}. These findings illustrate structural challenges that may discourage participation or retention among members of underrepresented groups.

Research has also examined the signals and environmental factors that influence developers' decisions to join or remain in software projects. For example, studies of open-source project participation have found that potential contributors evaluate factors such as team composition, communication style, and inclusivity when deciding whether to contribute to a project \cite{qiu2019signals}. In addition, broader research in diversity management emphasizes the role of organizational strategies, leadership commitment, and recruitment practices in shaping workforce diversity \cite{kreitz2008best}. Together, these studies suggest that improving diversity requires not only addressing structural barriers but also creating environments that actively support inclusion and equitable participation.

\section{Methodology}

\subsection{Matching Engineers with Teams}
The study examines how engineers are matched with teams within a large technology company and how this process influences team composition and diversity. At the time of the study, teams could recruit engineers through two main pathways: external hiring and internal transfers. External candidates typically applied through role-based job postings such as software engineering positions. After completing screening and interview stages, qualified candidates entered a team-matching phase where they were paired with specific teams that had open positions. During this process, hiring managers could evaluate multiple candidates simultaneously, while candidates could also interact with several teams before selecting a role. This matching process created a two-sided selection mechanism in which both teams and candidates influenced the final placement decision.

In addition to external hiring, existing employees within the organization could transfer between teams. Internal candidates were able to browse job listings on an internal platform and apply directly to teams that had open roles. Unlike external applicants, internal candidates often had access to additional contextual information about teams and hiring managers before applying, such as internal profiles, documentation, or examples of recent work. This visibility allowed employees to better evaluate team culture, leadership style, and technical focus when deciding whether to transfer. As a result, internal mobility functioned as another mechanism through which engineers could select teams that aligned with their professional goals and working preferences.

\subsection{Team Selection}

To study how diversity is cultivated within software engineering teams, the researchers first \textbf{identified candidate teams that could meaningfully represent diverse environments. Teams were defined as groups of engineers reporting directly to the same manager}. The researchers applied several criteria to narrow down the list of potential teams in order to ensure that the teams selected were both stable and organizationally meaningful units.
\begin{enumerate}
    \item One criterion involved team size: only teams consisting of five to fifteen engineers were considered. Smaller teams were excluded because diversity metrics would be unreliable with very few members, while very large teams were excluded because the manager-employee relationship might be less direct and the organizational structure less cohesive. Additionally, managers were required to have at least four years of management experience to ensure that they had sufficient experience managing teams across different organizational contexts and time periods.
    \item The researchers also applied demographic criteria to ensure that selected teams exhibited meaningful diversity. Teams were required to demonstrate diversity across both race or ethnicity and gender. For race and ethnicity, the researchers used a richness metric requiring at least four distinct racial or ethnic groups to be represented within a team. For gender diversity, the Gini-Simpson index was used to measure the probability that two randomly selected team members would have different genders. Teams were included only if their gender diversity index exceeded a specified threshold.
    \item In addition, at least one engineer from an underrepresented group needed to have remained on the team for more than two years and received a promotion during that period. This condition served as evidence that the team was capable not only of recruiting diverse engineers but also of supporting their retention and career advancement.
\end{enumerate}

\subsection{Interviewee Selection}

After identifying candidate teams, the researchers selected interview participants from two key roles within those teams: managers and software engineers. \textbf{Managers were prioritized if they had longer tenure within the organization and experience managing multiple teams, since such experience would allow them to compare diversity-related practices across different teams and contexts}. This perspective was important because managerial decisions and leadership behaviors often influence hiring practices, team culture, and opportunities for professional development within engineering teams. \textit{By interviewing managers with extensive organizational experience, the researchers aimed to gather insights into the strategies and decisions that contribute to building and sustaining diverse teams}.

For software engineers, the selection criteria emphasized individuals who belonged to one or more underrepresented groups, including women and individuals identifying with racial or ethnic groups other than White or Asian. Engineers were selected from different teams in order to capture a wide range of experiences rather than focusing deeply on a single team. Preference was also given to engineers who had spent longer periods within the organization or on their current team, since these participants could provide more detailed reflections on team culture and diversity practices over time. The researchers also aimed to ensure diversity within the interview sample itself by including participants from different racial, ethnic, and gender backgrounds. Additional steps were taken to include perspectives that might otherwise be underrepresented, such as engineers identifying as non-binary or belonging to demographic groups with very low representation in the workforce.

\subsection{Interviews}

The researchers conducted semi-structured interviews with selected participants in order to gather detailed insights about their experiences with diversity and inclusion within software engineering teams. All interviews were conducted remotely via teleconference and typically lasted approximately sixty minutes, although one interview lasted around thirty minutes. Semi-structured interviews were chosen to maintain consistency across participants while still allowing interviewees to discuss unexpected themes and experiences. Each interview began with an introduction to the study, followed by a request for permission to record the session and an explanation of how participants’ identities and responses would be protected to ensure confidentiality.

The interview questions were developed iteratively through several stages. Initially, the researchers created questions focused on recruitment and hiring practices related to diversity. These questions were later refined through feedback from colleagues and through pilot interviews with both managers and engineers. The final interview guides included separate sets of questions for managers and engineers. Managers were asked about their experiences building diverse teams, their hiring strategies, and how they fostered inclusive team cultures. Engineers were asked about their experiences working on different teams, their perceptions of diversity within those teams, and the factors influencing their decisions to join or leave teams. Open-ended questions were deliberately used so that participants could describe their experiences without being guided toward specific diversity practices discussed in existing literature.

\subsection{Data Analysis}

The collected interview data was analyzed using thematic analysis, a qualitative research technique that identifies patterns and themes across textual data. Interview recordings were first transcribed using automated tools and then manually reviewed to remove identifying information. The researchers conducted multiple rounds of open coding to categorize segments of text into conceptual codes, which were gradually refined and merged into a structured codebook. Throughout this process, the research team met regularly to review emerging codes and discuss interpretations of the data. The coding process continued until both code saturation and meaning saturation were reached, indicating that no new themes were emerging from the data. To improve the credibility of the findings, the researchers also performed member-checking by sharing summarized results with participants and allowing them to verify or comment on the interpretations of their responses.

\subsection{Self Disclosure}

The authors also acknowledged their positionality in relation to the study, following established practices in qualitative research. At the time the research was conducted, all authors were employees of the same organization from which the interview participants were recruited. Because of this relationship, the authors recognized that their professional context could potentially influence the research process and interpretation of findings. Although the researchers maintained autonomy in conducting and publishing the study, they acknowledged that their work was inevitably shaped by their organizational environment and their relationships with colleagues and leadership within the company.

\section{Results}

\subsection{Recruiting and Hiring}
\textbf{The first major results theme in the paper is that improving diversity in software engineering teams requires intentional action during recruiting and hiring, rather than treating representation as something determined only by the external pipeline}. A recurring finding was that managers on highly diverse teams did not regard lack of diversity as an unavoidable outcome; instead, they saw themselves as active agents who could influence team composition. The paper describes this as a shift toward self-empowering hiring managers. \textit{Participants explained that meaningful change began when managers first recognized that a diversity gap existed on their teams and then accepted responsibility for doing something about it}. In other words, the successful teams in the study were not simply lucky or passively diverse; their managers were described as people who reflected on their own incentives, assumptions, and prior practices, and then deliberately changed how they approached hiring. This is an important result because it frames diversity as something affected by everyday managerial choices rather than as a problem outsourced entirely to educational systems or upstream recruiting pipelines.

\textbf{A second major result is that these managers broadened the pool of candidates while simultaneously questioning the assumptions they brought into hiring decisions}. The paper shows that managers used several strategies to enlarge the candidate pool: \textit{they asked engineers from underrepresented groups to tap into their external networks, advertised roles more broadly, and built longer-term relationships through mentorship or outreach programs. At the same time, they reflected on what qualifications were genuinely necessary for success in the role. Instead of relying on a rigid or idealized image of the “right” engineer, they tried to distinguish core requirements from preferences that might unnecessarily narrow the field}. The findings suggest that this dual process—expanding where candidates come from while pruning back biased assumptions about who looks qualified—was central to more inclusive hiring. The study presents this as a practical counterweight to fast, convenience-driven hiring, where managers may be tempted to choose the first acceptable candidate rather than the best fit from a diverse pool.

The paper also emphasizes that inclusive hiring often required managers to slow the process down. \textbf{Several participants noted that managers took extra time to make sure job messages reached people from a broader range of backgrounds and to ensure those candidates had time to apply. This mattered because a rapid hiring process could easily favor the most immediately visible or conventionally credentialed candidates, which in practice often reproduced existing demographic patterns}. The study therefore presents time itself as a diversity resource: the willingness to wait, gather a broader candidate set, and resist settling for the first fit created better conditions for equitable hiring. What is notable here is that the paper does not portray these practices as symbolic gestures. Rather, interviewees described them as deliberate interventions that changed who was considered, how candidates were evaluated, and ultimately who joined the team.

\textbf{Another important result is that candidates from underrepresented groups were not only evaluating the technical role but also looking for evidence that managers were genuinely committed to inclusion}. The paper finds that an ongoing commitment to diverse candidates mattered well beyond the mechanics of screening and interviews. Managers signaled this commitment by investing in early-career training programs, learning about the experiences of engineers from underrepresented groups, participating in allyship or diversity-related efforts, and explicitly discussing how they supported career growth on their teams. Engineers reported using these signals to judge whether a team would be inclusive and whether they could succeed there. \textit{In some cases, engineers even examined publicly visible signs of a manager's involvement in diversity-related activities. This means recruiting was not a one-way process in which only managers assessed applicants; applicants from underrepresented groups were also actively assessing whether managers and teams were trustworthy, aware, and serious about inclusion}. The paper further argues that once managers built a visibly diverse and inclusive team, that itself became a positive signal that attracted more diverse candidates, creating a virtuous cycle.

\textbf{The final major result in this subsection is that individual manager effort was most effective when reinforced by organizational guidelines and accountability. Participants described leadership-level structures in which hiring managers had to consider multiple candidates and justify their choices to senior leaders such as vice presidents}. These accountability mechanisms made it harder to rely uncritically on intuition or bias and increased the likelihood that candidates from underrepresented groups would remain seriously considered throughout the hiring process. At the same time, the paper is careful not to treat top-down policy as sufficient on its own. Interviewees stressed that accountability measures worked best when combined with genuine bottom-up commitment from managers who believed in the goal of increasing diversity. \textit{The study also highlights a practical tension: equitable hiring often takes longer, but managers may feel pressure to fill roles quickly or risk losing headcount. As a result, the paper argues that leadership must do more than mandate inclusive processes; it must also protect the time and resources needed to carry them out}. This makes the recruiting and hiring findings especially strong, because they connect diversity not only to attitudes and intentions but also to institutional structures, incentives, and constraints.

\subsection{Technical Allyship}

\begin{definition}[Technical Allyship]
    Technical allyship refers to the deliberate actions taken by individuals within a technical team to support and advocate for engineers from underrepresented groups in ways that promote equity in technical work environments. It involves identifying and addressing structural or interpersonal biases that may undermine the recognition of these engineers’ technical contributions, while actively reinforcing confidence in their expertise and abilities. Unlike general expressions of support, technical allyship focuses specifically on behaviors within technical contexts—such as code reviews, task assignments, mentorship, and performance evaluations—that influence how engineers’ skills are perceived and valued.
\end{definition}

The second major theme identified in the study is the concept of technical allyship, which refers to actions taken by team members to support engineers from underrepresented groups in technical contexts. \textbf{The authors describe technical allyship as a form of allyship that goes beyond expressing support and instead involves concrete behaviors that address inequities in technical work environments. These actions help reinforce confidence in the technical skills of engineers from underrepresented groups and advocate for them when their expertise may be overlooked due to biases related to gender, race, or ethnicity}. The study highlights that technical allyship can be practiced at multiple levels within a team, including by managers, senior leaders, and coworkers.

Managers were found to play a particularly important role in practicing technical allyship. \textbf{Their decisions regarding performance evaluation, task allocation, and career development directly influence the technical opportunities available to engineers. Participants reported that effective managers were careful to gather sufficient evidence before evaluating performance and remained aware of biases that could lead to underestimating the abilities of engineers from underrepresented groups}. Managers also created stretch opportunities by assigning challenging technical work that allowed these engineers to demonstrate their capabilities and develop new skills. These practices helped counter stereotypes and ensured that engineers were not limited to less technical roles or overlooked for advancement opportunities.

\textbf{The study also highlights the importance of leadership-level allyship. Engineers from underrepresented groups often expressed concern about the lack of diversity among senior managers and technical leaders. Participants indicated that visible representation in leadership positions can signal that career advancement is achievable for individuals from diverse backgrounds}. Leaders who intentionally promote diversity within management ranks therefore play an important role in shaping organizational culture and long-term workforce diversity. Increasing representation at higher levels was described as having a cascading effect, influencing hiring practices and encouraging more inclusive environments across engineering teams.

\textbf{Finally, coworkers were also found to contribute significantly to technical allyship through everyday interactions. Engineers described how supportive colleagues helped build confidence by acknowledging contributions, engaging respectfully during code reviews, and offering constructive feedback during technical discussions. These interactions helped create an environment where engineers from underrepresented groups felt comfortable sharing ideas and participating in decision-making}. In addition, coworkers often provided mentorship, shared their own technical challenges, and advocated for colleagues’ contributions during evaluation or promotion processes. Such actions helped mitigate impostor syndrome and strengthened engineers’ sense of belonging within their teams.

\section{Discussion and Implications}
\subsection{Discussion}
The discussion section interprets the findings by emphasizing that the practices observed in highly diverse teams can potentially be transferred to other engineering teams seeking to improve diversity. The authors argue that diversity should not be treated as a fixed resource that teams compete for. \textit{A common argument is that diversity is a zero-sum situation, where increasing representation in one team necessarily reduces it elsewhere because the number of engineers from underrepresented groups is limited. The authors challenge this view, suggesting that such reasoning shifts responsibility away from teams and organizations. Instead, they argue that although the current talent pool may appear limited at any given time, it is influenced by broader systemic factors and can expand over time through cultural and structural change}. Engineering teams therefore share responsibility for improving inclusivity within their own environments rather than attributing low diversity solely to upstream pipeline issues.

The authors also discuss practical challenges associated with implementing the practices identified in the study. \textit{Although the practices appear effective in diverse teams, they may involve trade-offs that make adoption difficult in other organizational contexts. For example, inclusive hiring practices often require managers to spend more time identifying and evaluating candidates from diverse backgrounds. However, organizations frequently operate under pressure to fill open roles quickly, which can discourage managers from taking the additional time required to broaden candidate pools}. Additionally, the study found that the managers who successfully cultivated diverse teams were strongly committed to diversity goals and felt personal responsibility for improving representation. The authors suggest that managers who lack this level of commitment or who follow diversity initiatives only superficially may be less successful in achieving meaningful change.

Another important point raised in the discussion concerns the relationship between diversity practices and broader management behaviors. \textit{Many of the practices identified in the study overlap with characteristics of effective engineering leadership described in prior research, such as supporting engineers' growth, encouraging experimentation, and fostering respectful collaboration. At the same time, the study highlights additional behaviors that are particularly important for inclusion, such as carefully evaluating performance to avoid bias and encouraging a culture of mutual respect and mentorship within teams}. The authors also emphasize the role of coworkers, noting that inclusive environments are not created solely through management decisions but also through everyday interactions among engineers. Actions such as recognizing colleagues' contributions, providing constructive feedback, and advocating for fair evaluation help reinforce inclusive team cultures and support the technical confidence and career progression of engineers from underrepresented groups.

\subsection{Implications}
\subsubsection{Engineers}

Engineers play an important role in shaping inclusive team environments through their everyday interactions and collaboration practices. Although diversity initiatives are often discussed at the managerial or organizational level, the study highlights that individual engineers can actively contribute to more equitable and supportive technical environments.

\begin{itemize}
    \item Act as technical allies by reinforcing confidence in colleagues from underrepresented groups and acknowledging their technical contributions during discussions, code reviews, and presentations.
    \item Advocate for fairness in team processes by speaking up when disparities in recognition, workload, or opportunity distribution are observed.
    \item Avoid reinforcing stereotypes or informal role assignments that may push certain engineers into less technical responsibilities.
    \item Encourage open collaboration by providing constructive feedback and fostering respectful technical discussions that value diverse perspectives.
    \item Request transparency from management regarding hiring, promotion, and evaluation processes to ensure equitable treatment within teams.
\end{itemize}

\subsubsection{Managers}

Managers have significant influence over team composition, culture, and career development opportunities. The study emphasizes that managers must actively examine their own practices and decisions to ensure they are supporting diversity and inclusion within their teams.

\begin{itemize}
    \item Evaluate diversity and inclusion practices within the team and reflect critically on how hiring, evaluation, and task assignment decisions are made.
    \item Develop awareness and empathy regarding the experiences of engineers from underrepresented groups by engaging with training programs, allyship initiatives, and educational resources.
    \item Seek regular informal feedback from team members to understand their experiences and improve team culture.
    \item Implement inclusive hiring strategies by expanding candidate pools and examining potential biases when evaluating applicants.
    \item Carefully assess employee performance by gathering comprehensive evidence rather than relying on intuition or limited impressions.
    \item Provide stretch opportunities and challenging technical tasks to engineers from underrepresented groups to support their growth and visibility.
    \item Clearly communicate commitment to career development and mentorship when interacting with potential candidates during recruiting.
\end{itemize}

\subsubsection{Leaders}

Senior leaders and organizational decision-makers play a crucial role in shaping policies and structural mechanisms that influence diversity outcomes across teams. Their actions can reinforce inclusive practices and ensure accountability for diversity initiatives throughout the organization.

\begin{itemize}
    \item Increase diversity in leadership positions to provide visible representation and demonstrate that career advancement is achievable for engineers from underrepresented groups.
    \item Establish structured hiring and promotion practices that require managers to consider diverse candidate pools during recruitment.
    \item Implement accountability mechanisms that require managers to document and justify hiring decisions, helping reduce the influence of bias.
    \item Recognize and reward managers who actively improve diversity and foster inclusive team cultures.
    \item Provide managers with sufficient time and organizational support to conduct inclusive hiring processes without the pressure to fill positions rapidly.
\end{itemize}

\subsection{Limitations}
The study has several limitations that affect the generalizability and interpretation of its findings. First, the research was conducted within a single large technology company, and all interview participants were employees working in the United States. As a result, the organizational context, hiring practices, and team structures examined in the study may differ significantly from those in smaller companies, startups, or open-source communities. For example, organizations with fewer resources or less formal recruitment processes may not be able to implement some of the practices identified in the study, such as extensive hiring pipelines or structured accountability mechanisms for recruitment decisions. Consequently, the transferability of the findings may be limited to organizations with similar scale, infrastructure, and workforce demographics \cite{easterbrook2008selecting}.

Another limitation arises from the scope of diversity considered in the research and the potential for bias in qualitative self-report data. The study primarily focuses on diversity in terms of race, ethnicity, and gender using broad demographic categories. Other dimensions of diversity—such as socioeconomic background, disability, or cultural identity—were not examined and may require different strategies to address. Additionally, because the study relied on interviews, participants may have been influenced by social desirability bias, potentially presenting their teams or practices in a more positive light than reality. Finally, the study may also be subject to survivorship bias, since it interviewed current employees on successful teams rather than engineers who left the organization or the profession due to negative experiences \cite{nederhof1985methods}.